\documentclass[12pt,a4paper]{article}
\usepackage[utf8]{inputenc}
\usepackage[T1]{fontenc}
\usepackage[english,indonesian]{babel}
\usepackage{amsmath}
\usepackage{amsfonts}
\usepackage{amssymb}
\usepackage{graphicx}
\usepackage{longtable}
\usepackage{hyperref}
\usepackage{geometry}
\geometry{margin=1in}

\title{\textbf{Judul Berbahasa Indonesia} \\ \emph{Title in English}}
\author{Nama Penulis 1$^{1}$, Nama Penulis 2$^{2}$ \\ $^{1}$Afiliasi 1, Alamat 1, surel1@email.com \\ $^{2}$Afiliasi 2, Alamat 2, surel2@email.com}
\date{}

\begin{document}

\maketitle

\begin{abstract}
\textbf{Abstract} (Bahasa Inggris)

The abstract should addresses the following: key aspects of the literature review, problem under investigation or research questioning, clearly stated hypothesis or hypotheses, method used (including brief description of the study design, sample, and sample size), study results, implications (i.e. why the study is important, applications of the results or findings). (Times New Roman, Size 11)

Keywords: ..........., ..........., ..........., ..........., ...........,

\textbf{Abstrak} (Bahasa Indonesia)

Abstrak memuat: aspek kunci dari review literatur, pertanyaan penelitian, hipotesis, metode pemecahan masalah, hasil analisis, rekomendasi terhadap hasil temuan. (Times New Roman, Size 11)

Kata Kunci: ..........., ..........., ..........., ..........., ...........,

(Abstrak paling banyak 150 kata dalam Bahasa Indonesia dan 200 kata dalam Bahasa Inggris. Abstrak dilengkapi dengan kata kunci antara 3 sampai dengan 5 kata/frasa.)
\end{abstract}

\section{\textbf{PENDAHULUAN}}

\begin{enumerate}
\def\labelenumi{\Alph{enumi}.}
\item \textbf{Latar Belakang}

Latar belakang memuat penjelasan mengapa topik diangkat sebagai tulisan dan rumusan masalah yang diidentifikasi.

\item \textbf{Pertanyaan Penelitian}

Menjelaskan masalah yang muncul dalam penyusunan artikel jurnal ilmiah.

\item \textbf{Tujuan}

Menjelaskan tujuan penelitian yang dilakukan.

\item \textbf{Manfaat}

Menjelaskan manfaat dari penelitian yang dilakukan.
\end{enumerate}

\section{\textbf{METODE}}

Jenis penelitian, teknik dan cara pengumpulan data, dan metode analisis.

\section{\textbf{HASIL DAN BAHASAN}}

Analisis dan bahasan terhadap data yang diperoleh. Subjudul bahasan menggunakan huruf kapital (A, B, C)

\subsection{\textbf{Subbahasan I}}

Berikut contoh penggunaan tabel dengan tata letak sesuai lebar kolom

\begin{table}[h]
\centering
\caption{Jumlah Penulisan Jurnal Kearsipan}
\begin{tabular}{|l|l|l|l|}
\hline
No. & Tahun & Jumlah & Keterangan \\
\hline
1 & 2019 & 14 & - \\
2 & 2020 & 12 & - \\
3 & 2021 & 15 & - \\
\hline
\end{tabular}
\label{tab:example1}
\end{table}

Sumber: Direktorat Sistem Kearsipan

Berikut contoh penggunaan gambar dengan tata letak sesuai lebar kolom

\begin{figure}[h]
\centering
\includegraphics[width=0.5\textwidth]{media/image1.jpg}
\caption{Rak Simpan Arsip}
\label{fig:example1}
\end{figure}

Sumber: google (daring menyertakan waktu akses)

\subsection{\textbf{Subbahasan II}}

Berikut contoh tabel dengan tata letak sesuai lebar halaman

\begin{table}[h]
\centering
\caption{Instansi Asal Penulis Jurnal Kearsipan}
\begin{tabular}{|l|l|l|l|}
\hline
No. & Tahun & Tahun & Keterangan \\
\hline
1 & Arsip Nasional Republik Indonesia & 2020 & - \\
2 & Kementerian ABCD & 2021 & - \\
3 & Lembaga Pemerintah Non Kementerian & 2022 & - \\
\hline
\end{tabular}
\label{tab:example2}
\end{table}

Sumber: Direktorat Sistem Kearsipan

\subsection{\textbf{Subbahasan III, dst}}

Berikut contoh penggunaan gambar dengan tata letak sesuai lebar halaman

\begin{figure}[h]
\centering
\includegraphics[width=0.8\textwidth]{media/image2.jpg}
\caption{Instansi Asal Penulis Jurnal Kearsipan}
\label{fig:example2}
\end{figure}

Sumber: google (daring menyertakan waktu akses)

\section{\textbf{SIMPULAN}}

\section{\textbf{REKOMENDASI}}

\section{\textbf{PERNYATAAN TERIMA KASIH}}

\section{\textbf{DAFTAR PUSTAKA}}

Mencakup kutipan dalam bentuk catatan dalam atau catatan kaki. Semua sumber kutipan harus dimuat dalam daftar sumber, berlaku pula sebaliknya. Dapat menggunakan aplikasi Mendeley atau Zotero dan gaya kutipan American Sociological Association (ASA) Style Guide. Sila merujuk https://www.utm.utoronto.ca/sociology/sites/files/sociology/documents/2023-08/ASA_Style_Guide\%207th\%20ed.pdf atau https://www.apsu.edu/writingcenter/writing-resources/ASA-Style-7th-Edition-2024pdf.pdf

\textbf{Peraturan Perundang-Undangan}

Indonesia, R. 2009. Undang-Undang No. 43 Tahun 2009 tentang Kearsipan. Lembaran Negara RI Tahun 2009. Jakarta: Sekretariat Negara.

\textbf{Buku}

Rustam, M. 2014. Modul I: Pengantar Pengelolaan Arsip Elektronik. Jakarta: Universitas Terbuka.

Smallwood, R. F. 2013. Managing Electronic Records. Methods, Best Practices, and Technologies. New Jersey: John Wiley\& Sons, Inc.

\textbf{Artikel}

\textbf{dalam bunga rampai}

Pratama, Raistiwar. 2016. ``Membaca Pram dan Orwell,'' dalam Nadia Fauziah Dwiandari. Arsip dan Budaya Indonesia. Jakarta: Ikatan Arsiparis Arsip Nasional Republik Indonesia.

\textbf{dalam jurnal}

Pratama, Raistiwar. 2022. ``Arsip adalah Tempat: Definisi, Paradigma, dan Kepemilikan,'' Diplomatika 5 (2): 72 -- 91.

\textbf{Laman dalam Jaringan}

Campbell-Miller, Jill; and Ryan Kirkby. 2018. ``Historian, Meet Archivist: Researching the History of Complex Organizations.'' http://activehistory.ca/2018/07/historian meet-archivist-researching-the-history-of complex-organizations/ (Accessed July 4th 2018).

\textbf{Sumber Lisan Wawancara.}

Nama terwawancara/Nama samaran/ Inisial. Tempat. Waktu.

\textbf{Tugas Akhir, Skripsi, Tesis, Disertasi}

Pratama, Raistiwar. 2017. ``Archives and the Archipelago: The Influence of Dutch Archivistiek on Indonesian Archival Practices.'' Unpublished Master Thesis. University of Leiden.

\end{document}
